\documentclass[10pt]{article}

\usepackage[margin=1in]{geometry}
\usepackage{booktabs}
\usepackage{amsmath}
\usepackage{amssymb}
\usepackage{graphicx}
\usepackage{hyperref}
\usepackage{xcolor}

\title{Semantic-Scale Support Mismatch in Remote-Sensing Object Detection}
\author{Anonymous Authors}
\date{Draft: 2026-06-21, v0.9 stricter devil-reviewer Gaussian support-field contract}

\newcommand{\sise}{\mathrm{SISE}}
\newcommand{\logarea}{\log \mathrm{area}}

\begin{document}
\maketitle

\begin{abstract}
Remote-sensing object detectors operate in scenes where object categories have
strong physical scale regularities.  We study a reliability failure in which a
detector localizes an object but assigns a high-confidence class whose predicted
box scale lies in a low-support region for that class.  We call this failure
semantic-scale support mismatch.  The stricter framing is not a heuristic tail
penalty; it is a Gaussian semantic-scale support field.  A semantic confidence
claim should be compatible with a class- and domain-conditioned generative law
over object log-area, measured as a likelihood-ratio support score against the
domain background scale distribution.  The currently validated implementation
is the source-disjoint single-Gaussian case, which yields Scale-Inconsistent
Semantic Error (SISE): a localized wrong-label detection whose predicted class
has low support under the predicted scale.
Across open-vocabulary and closed-set remote-sensing settings, a continuous
Gaussian support-risk feature predicts pair-level high-confidence risk.
Current audits
show 6/6 datasets passing the support-law gate, OVD and closed-set domain
correlations of 0.823534 and 0.636249, and DOTA2 high-risk pair recurrence of
1.000 across detector families.  We introduce G-S3C, a Gaussian semantic-scale
support calibration family that modulates logits or scores using
$p(\logarea\mid \mathrm{class})$ without using Gaussian approximations to IoU.
The stricter method object is not more suppression; it is an AP-aware
compatibility contract between the semantic score surface, the Gaussian support
field, and the ranking or assignment surface that produces detections.  G-S3C
reduces high-confidence SISE with AP non-regression in selected settings.
Against score-only global Platt, classwise Platt, and isotonic calibration,
G-S3C/G2 reduces fixed-correct SISE more on two of three closed-set audits with
measurable queue risk.
The current evidence also defines clear boundaries: E-P2-clean does not pass a
causal path gate, G3-v1 trainable consistency does not beat matched no-G3
controls, and the completed G3-v2 matched run is only marginally mAP-positive
while failing the AP-risk queue gate.  The current paper is therefore a
diagnostic and deployable reliability study.  The 9.5 method target is
BASS-GSF: a Bayesian semantic-scale support detector using a Gaussian support
field, feature-conditioned posterior compatibility, semantic-scale virtual
outliers, and assignment/ranking coupling.  This is the only route by which the
current work can honestly claim method novelty at the stricter 9.5 level.
The completed AP-aware RankDelta/P2 closure does not validate that route.  A
first two-GPU BASS-GSF-Lite pilot improves HRRSD density e2 from mAP 0.860457
to 0.861822 and log-z SISE from 143 to 66, but remains below the no-G3 e2 AP
control 0.862878.  The best P1 RankDelta variant is a useful pilot
(mAP 0.862488, localized wrong 6194, risk-weighted log-z 12.0345), yet still
misses the strict no-G3 AP/top-$K$ gate.  P2 assign-rank lowers localized wrong
to 6122 but drops mAP to 0.857772, while P2 posterior-rank reaches mAP
0.860845 but worsens localized wrong to 6720 and risk-weighted log-z to
16.0802.  P3A then implements the resulting AP-constrained Gaussian semantic
support projection as a validation-set oracle: correct positives are protected
and only low-support localized wrong-class detections are demoted.  Against
the no-G3 e2 control, the selected P3A z=2.0, score-floor 0.1 row improves mAP
from 0.862878 to 0.862961, top5000 precision from 0.7168 to 0.7170, and
top5000 log-z SISE from 48 to 0, with changed-correct positives equal to 0 and
TP rank harm equal to 0.  P3D transfers the same AP-constrained projection to
DIOR-R: mAP improves from 0.644687 to 0.645981, top5000 precision remains
1.0000, all-risk log-z false alarm falls from 2236.0471 to 1827.5576,
score>=0.5 risk-weighted log-z falls from 6.1242 to 0.0037, 54 AP-sensitive
SISE false positives are removed, and both changed-correct positives and TP
rank harm remain 0.  This closes the dataset-transfer breadth gate for the
8/10 machine audit.  A non-oracle P4/GSRL score-only ranker using only
predicted class, predicted area, score, and train-set Gaussian priors reaches
only mAP 0.644693 on DIOR-R, below P3D's 0.645981, and fails rank-tail safety.
Thus non-oracle controls and detector-family transfer remain open, and
Gaussian support must be coupled to AP-aware ranking or assignment rather than
used as a free-standing score prior.
The next extension raises the support variable from scalar scale to geometric
support without becoming a Gaussian box-loss paper.  P5B represents an object
by $g=[\log(\mathrm{area}),\log(\mathrm{long\ side}/\mathrm{short\ side})]$
and models $p(g\mid c,d)$ as a class/domain-conditioned Gaussian prior over
semantic geometry.  Multi-dataset geometry priors have been materialized for
HRRSD, ShipRS, DIOR-R, and xView, and the implementation now exposes an
optional head-level geometry log-prob gate for train-time hard-negative and
positive selection.  The completed P5B runs are negative: the mild geometry
gate reaches mAP 0.8583 and the stronger gate reaches 0.8550, both below the
no-G3 e2 control 0.862878, while rank-tail audits show hundreds of
correct-positive rank drops.  This failure is informative: Gaussian geometry
support is useful only if it is coupled to AP/rank protection, not as an
unconstrained delta loss.
\end{abstract}

\section{Introduction}

Remote-sensing object detection is shaped by severe scale variation, arbitrary
orientation, dense object layouts, and large viewpoint changes.  DOTA and
DOTA-v2 make these challenges explicit at benchmark scale
\cite{xia2017dota,ding2021dotav2}.  Modern detectors address scale through
feature pyramids, scale-aware assignment, receptive-field specialization, and
distributional localization \cite{lin2016fpn,li2019tridentnet,tian2019fcos,zhang2019atss,li2020gfl}.
Remote-sensing detectors additionally model orientation and context through
rotation-aware architectures and specialized backbones
\cite{han2020s2anet,han2021redet,yu2023h2rboxv2,li2024lsknet}.
Open-vocabulary detectors add language-conditioned class scores
\cite{gu2021vild,li2021glip,minderer2022owlvit,liu2023groundingdino,zhou2022detic,cheng2024yoloworld,pan2024laedino,huang2025openrsd}.

These lines do not directly answer a deployment reliability question:

\begin{quote}
When a detector confidently predicts a class label for a localized object, is
the predicted object scale plausible for that class?
\end{quote}

Our thesis is that a remote-sensing class prediction is not only a semantic
classification event; it is also a physical support claim.  If a detector says
``harbor'', ``airplane'', or ``small vehicle'', the predicted box scale should
lie in a plausible region of the class/domain scale distribution.  A
high-confidence semantic logit outside that support is a reliability violation
even when the box is localizable and even when aggregate AP hides the failure.
This turns scale from a feature-pyramid engineering variable into a
probabilistic compatibility variable between semantics and the physical world.

This question is not equivalent to localization quality.  A box can match a
true object while the predicted class is physically implausible for the object's
scale.  We call this semantic-scale support mismatch.  The core diagnostic is
class-conditioned support over log area:
\begin{equation}
  a_i = \log(\mathrm{area}(b_i)), \qquad
  p(a\mid c) = \mathcal{N}(\mu_c,\sigma_c^2), \qquad
  z_{i,c} = \frac{a_i-\mu_c}{\max(\sigma_c,\epsilon)}.
\end{equation}

For a localized prediction with predicted class $c_i$ and matched ground-truth
class $g_i$, SISE is counted when the prediction is wrong and the predicted
class has low scale support:
\begin{equation}
  \sise_{\tau} =
  \mathbb{1}[\mathrm{IoU}(b_i,b_i^{gt}) \ge \tau_{\mathrm{iou}}]
  \mathbb{1}[c_i \ne g_i]
  \mathbb{1}[|z_{i,c_i}| \ge \tau_z]
  \mathbb{1}[s_{i,c_i} \ge \tau_s].
\end{equation}

\paragraph{Gaussian support-field view.}
The higher-level claim is a reliability law, not merely a penalty rule.
Semantic confidence should be supported by a generative scale field:
\begin{equation}
  p_\phi(a\mid c,d)=\sum_k \pi_{c,d,k}
  \mathcal{N}(a;\mu_{c,d,k},\sigma^2_{c,d,k}), \qquad
  S_c(a,d)=\log p_\phi(a\mid c,d)-\log p_0(a\mid d).
\end{equation}
Here $d$ denotes the domain and $p_0(a\mid d)$ is the background object-scale
support in that domain.  The current G-S3C evidence uses the minimal $K=1$
class-conditioned Gaussian distribution.  The likelihood-ratio form matters: a class should
not be punished merely for having a rare absolute size if that size is rare for
all classes in the domain, and it should be punished more strongly when the
same size is common in the domain but low-support for the asserted class.  The
9.5 method route is to make this support field trainable through BASS-GSF
rather than stopping at post-hoc Gaussian calibration.

The Gaussian distribution is not decoration.  It is the smallest auditable
support law that turns a scale heuristic into a probabilistic semantic claim:
each class obtains a mean, uncertainty, tail probability, and likelihood ratio
against the domain background.  Mixture and hierarchical Gaussians are then
principled extensions for multi-modal classes, rare classes, and domain shift.
The resulting problem is higher than size-based calibration: a detector should
make class claims that are compatible with a physical support distribution.
This also raises the proof burden.  A thresholded $|z|$ rule can justify a
post-hoc filter, but a Gaussian support-distribution claim creates three
reviewable objects: the empirical class/domain scale law, the likelihood-ratio
support score, and the train-time posterior or ranking mechanism that must keep
AP intact.

We call this stronger requirement the Gaussian semantic support contract.  For
every high-confidence localized class claim, the detector must provide a score
that is simultaneously semantically high, scale-supported under
$p_\phi(a\mid c,d)$, and rank-safe under AP-sensitive evaluation.  This
contract is why the negative BASS-delta evidence matters: a method can reduce
low-support false semantics while still damaging the ranking surface.  That is
not a 9.5 method result; it is a diagnosis of what the next method must
protect.

The broader Gaussian application is therefore not limited to one scalar
log-area prior.  It gives a common probabilistic language for five detector
objects: class/domain scale support $p_\phi(a\mid c,d)$, candidate posterior
compatibility $q_\theta(a\mid h,c)$, low-likelihood semantic-scale virtual
outliers, AP-constrained score projection, and deployment risk sets.  This
connects remote-sensing scale regularity to probabilistic detection,
energy-based/OOD modeling, assignment, and risk-control literature while
keeping the scope precise: the paper studies semantic support of class claims,
not Gaussian approximations to IoU.

The same support-field language can encode geometric prior knowledge.  Many
remote-sensing targets are geometrically ordered, approximately symmetric, and
smoothly bounded under oriented-box annotation.  We therefore introduce a
geometry support coordinate
\begin{equation}
  g_i=\left[\log(\mathrm{area}(b_i)),\,
  \log\frac{\mathrm{long\ side}(b_i)}{\mathrm{short\ side}(b_i)}\right].
\end{equation}
The second coordinate is invariant to rotation and quadrilateral point order in
DOTA-style annotations.  A class-conditioned $p_\phi(g\mid c,d)$ can distinguish
near-symmetric categories from elongated or ordered categories without
measuring similarity between two boxes.  This is the route by which Gaussian
modeling enters backbone/encoder/adapter/head decisions as semantic support,
not as GWD, KLD, NWD, or an IoU surrogate.

\paragraph{Current evidence status.}
This draft is intentionally evidence-safe.  The support-risk law is a
predictive reliability claim, deployment practicality is positive, and
literature positioning is complete enough for review.  However, the
individual-level causal path gate and network-internal BASS-GSF method gate
remain open.

\section{Related Work and Closest Competitors}

\paragraph{Scale-aware detection.}
FPN, TridentNet, FCOS, ATSS, and GFL address feature scale, receptive field,
dense assignment, and distributional localization
\cite{lin2016fpn,li2019tridentnet,tian2019fcos,zhang2019atss,li2020gfl}.
They improve detection across object sizes, but they do not explicitly model
whether a predicted class is plausible under class-conditioned physical scale
support.

\paragraph{Remote-sensing detection.}
DOTA/DOTA-v2 and oriented detectors such as S2A-Net, ReDet, H2RBox-v2, and
LSKNet motivate the remote-sensing setting and provide detector families where
semantic-scale mismatch recurs
\cite{xia2017dota,ding2021dotav2,han2020s2anet,han2021redet,yu2023h2rboxv2,li2024lsknet}.

\paragraph{Open-vocabulary detection.}
CLIP, ViLD, GLIP, OWL-ViT, Grounding DINO, Detic, YOLO-World, LAE-DINO, and
OpenRSD expand semantic coverage
\cite{radford2021clip,gu2021vild,li2021glip,minderer2022owlvit,liu2023groundingdino,zhou2022detic,cheng2024yoloworld,pan2024laedino,huang2025openrsd}.
Language alignment can produce broad semantic scores, but it does not guarantee
scale-plausible class confidence in remote-sensing scenes.

\paragraph{Calibration, probabilistic detection, and OOD.}
Detector calibration work shows that AP and calibration must be evaluated
jointly, and that cheap post-hoc baselines can be strong
\cite{kuzucu2024detcalib}.  Cal-DETR uses uncertainty-guided logit modulation
\cite{munir2023caldetr}.  Probabilistic object detection evaluates semantic and
spatial uncertainty jointly \cite{hall2018probod}.  VOS synthesizes
low-likelihood feature outliers for OOD uncertainty \cite{du2022vos}.  G-S3C
differs by modeling a specific reliability axis: $p(\logarea\mid c)$.  The
current package therefore includes score-only Platt and isotonic baselines that
can use score/correctness labels but cannot read log-area surprise or
class-scale priors.  BASS-GSF borrows the outlier idea only after changing the
low-support variable: a virtual negative is a class-scale claim that is unlikely
under $p_\phi(a\mid c,d)$ but plausible under the domain background
$p_0(a\mid d)$.

\paragraph{Gaussian uncertainty and energy models.}
Several detectors already use Gaussian modeling for localization uncertainty,
including Gaussian YOLOv3 and uncertainty-aware dense detection
\cite{choi2019gaussianyolo,park2021uadet}.  These works justify Gaussian
uncertainty as a detector-internal object, but their random variable is box
localization, not semantic-scale support.  Deep feature OOD work uses
class-conditional Gaussian models and Mahalanobis scores, while energy-based
OOD detection reframes confidence as an energy surface rather than raw softmax
\cite{lee2018mahalanobis,liu2020energyood}.  This is the closest conceptual
support for our larger view: $S_c(a,d)$ should be read as a semantic-scale
energy or likelihood ratio, and BASS-GSF should shape the detector decision
surface rather than only threshold low-support predictions.

\paragraph{Assignment, ranking, and risk control.}
PAA, OTA, and TOOD show that detector performance depends on how training
samples and task-aligned scores are assigned, not only on the final classifier
\cite{kim2020paa,ge2021ota,feng2021tood}.  This literature explains why our
P2 results fail: Gaussian support can reduce some wrong labels, but without
AP-critical assignment and ranking constraints it can still damage the
ordering surface.  Conformal object-detection work and conformal risk control
provide an orthogonal deployment view: reliable detection may require
prediction sets or risk bounds rather than a single overconfident class
\cite{andeol2023conformalrisk}.  We use this as motivation for support-aware
risk gates, not as a substitute for AP-risk detector validation.

\paragraph{Gaussian box methods.}
GWD, KLD, and NWD model boxes as Gaussian distributions for rotated or tiny
object localization \cite{yang2021gwd,yang2021kld,wang2021nwd}.  G-S3C does
not use Gaussian distributions for box overlap or IoU approximation.  This
boundary is non-negotiable because the Gaussian-box route is already crowded.
Gaussian is used only for class-conditioned semantic-scale support: whether a
semantic claim is plausible under the physical log-area support of that class.
A Gaussian box loss asks whether two boxes match geometrically; a Gaussian
semantic-scale support field asks whether the class claim is physically
supported.

\section{Method}

\subsection{G-S3C}

G-S3C uses class-conditioned log-area priors to modulate class logits or scores.
The safe L0 guard preserves the existing S3C behavior as a fallback.  G1 uses a
continuous Gaussian logit energy:
\begin{equation}
  \Delta \ell_{i,c} = -\beta\,\mathrm{softplus}(|z_{i,c}| - z_0),
  \qquad
  \ell'_{i,c} = \ell_{i,c} + \Delta \ell_{i,c}.
\end{equation}
The default parameters are $z_0=4.0$ and $\beta=1.38629436112$.

G2 extends this idea with a lightweight scale-aware logit adapter using
$|z|$, signed $z$, Gaussian log probability, log standard deviation, and a
valid-prior mask.  By default, the adapter is constrained to non-positive
deltas to avoid creating new high-confidence false positives.

\subsection{Gaussian Support Field and BASS-GSF Gate}

BASS-GSF is the planned network-internal method.  It should model compatibility
between dense features, class semantics, and class/domain scale support:
\begin{equation}
  q_\theta(a\mid h_i,c)=
  \mathcal{N}(m_\theta(h_i,c),s_\theta^2(h_i,c)), \qquad
  S_{i,c}=\log p_\phi(a_i\mid c,d)-\log p_0(a_i\mid d).
\end{equation}
The intended objective couples detection loss, posterior compatibility,
semantic-scale negatives, and ranking:
\begin{equation}
\begin{aligned}
  \mathcal{L}_{\mathrm{BASS}} =
  \mathcal{L}_{\mathrm{det}}
  &+ \lambda_{\mathrm{pos}}
  \mathrm{KL}\!\left(q_\theta(a\mid h_i,c_i)\,\|\,p_\phi(a\mid c_i,d)\right)\\
  &+ \lambda_{\mathrm{neg}}
  \mathbb{E}_{c^-}\left[\max(0,\tau+S_{i,c^-}-S_{i,c_i})\right]
  + \lambda_{\mathrm{rank}}\mathcal{L}_{\mathrm{rank}}(S,\mathrm{score}).
\end{aligned}
\end{equation}

\paragraph{Gaussian likelihood-ratio decision field.}
The stricter formulation treats $S_{i,c}$ as a log-likelihood ratio, not as a
tail penalty:
\begin{equation}
  G_{i,c} =
  \log p_\phi(a_i\mid c,d) - \log p_0(a_i\mid d), \qquad
  F_{i,c} = \ell_{i,c} + \eta G_{i,c}
  + \rho \log q_\theta(a_i\mid h_i,c).
\end{equation}
Here $G_{i,c}$ asks whether the class claim is supported relative to the
domain's object-scale background, and $F_{i,c}$ is the detector decision field
after semantic evidence and Gaussian support evidence are combined.  This is
the conceptual lift: a detector class score becomes a posterior-compatible
support claim under a Gaussian generative field.  The random variable is
semantic log-area support, not a Gaussian box representation; no term in
BASS-GSF approximates bbox IoU, GWD, KLD, or NWD.

\paragraph{Geometry-aware Gaussian support.}
P5B extends the random variable from scalar log-area to semantic geometry:
\begin{equation}
  g_i=\left[\log(\mathrm{area}(b_i)),
  \log\frac{\mathrm{long\ side}(b_i)}{\mathrm{short\ side}(b_i)}\right],
  \qquad
  G_{i,c}^{g}=\log p_\phi(g_i\mid c,d)-\log p_0(g_i\mid d).
\end{equation}
The second coordinate injects the remote-sensing prior that many objects have
stable aspect structure: symmetric classes concentrate near zero log-aspect,
whereas ships, bridges, runways, and harbors occupy elongated support regions.
Right-angle error is retained only as an annotation-quality and boundary
smoothness audit.  The trainable support signal should come from log-area and
log-aspect likelihood, keeping the Gaussian model in semantic support space
rather than in box-distance space.

\paragraph{AP-constrained support projection.}
The P0--P2 boundary results show that reducing unsupported hard negatives is
not enough.  The trainable method must project the detector score field onto a
Gaussian-compatible region while preserving AP-critical orderings:
\begin{equation}
\begin{aligned}
  F^\star = \arg\min_{\widetilde{F}}\;
  &\sum_{i,c}(\widetilde{F}_{i,c}-F_{i,c})^2
  + \lambda_G \mathcal{R}_G(\widetilde{F};G)\\
  \mathrm{s.t.}\quad
  &\widetilde{F}_{i,g_i}-\widetilde{F}_{j,c^-}
  \ge
  F_{i,g_i}-F_{j,c^-}-\epsilon,
  \quad (i,j,c^-)\in\mathcal{C}_{\mathrm{AP}} .
\end{aligned}
\end{equation}
$\mathcal{C}_{\mathrm{AP}}$ denotes positive--hard-negative pairs that affect
top-$K$ ranking, and $\mathcal{R}_G$ penalizes decisions that assign high
scores to classes with low Gaussian support likelihood ratio.  This makes the
next BASS-GSF route an AP-safe Gaussian support projection rather than a larger
negative delta.

The negative classes $c^-$ should include virtual semantic-scale outliers:
class claims whose scale is unlikely under $p_\phi(a\mid c,d)$ but plausible
under $p_0(a\mid d)$.  This is the method route that could lift the paper from
validated reliability calibration to a stronger trainable detector.

The strict design rule is that $S_{i,c}$ must influence the detector at a
decision point where ranking can still change and correct positives are
explicitly protected.  If $S_{i,c}$ is used only as an offline audit, the
contribution remains diagnostic.  If it only subtracts logits from hard
negatives, the completed train-time delta run already shows the failure mode:
SISE drops while mAP and top5000 precision can still fall.  The 9.5 version
must therefore couple support compatibility to at least one AP-sensitive
mechanism: assignment weights, pairwise top-$K$ ranking, or feature-conditioned
posterior density with a positive keep term.

For this route to count as a real method rather than a restated calibration
rule, BASS-GSF must satisfy four reviewer-facing tests.  First, the Gaussian
field must enter training or pre-topK ranking, not only post-hoc reporting.
Second, posterior compatibility must be class-conditioned and feature-aware, so
the model learns when a candidate is a plausible instance of a class at a
specific scale.  Third, semantic-scale virtual outliers must create hard
negative pressure in regions where ordinary assignment supplies too few
examples.  Fourth, the resulting detector must improve an AP-risk Pareto gate
against matched no-BASS controls; a lower SISE count alone is insufficient.

The stricter 9.5 contract is falsifiable: under matched no-BASS controls with
the same seed, schedule, backbone, and dataset split, a BASS-GSF variant is a
valid main method only if it preserves or improves mAP and top-$K$ precision
while reducing low-support semantic risk.  A single Gaussian prior can lift
problem anatomy and deployment reliability, but it cannot by itself lift method
novelty, effectiveness, applicability breadth, and ICLR readiness to 9.5.  The
Gaussian support field must change the detector decision surface without
breaking AP-sensitive ranking.

BASS-GSF-Lite implements the first deployable part of this route.  Given a
trained density head, it converts feature-conditioned posterior log probability
into a non-positive pre-topK logit energy:
\begin{equation}
  \Delta \ell_{i,c} =
  -\beta\,\mathrm{softplus}(\tau_{\log p} -
  \log p_\theta(a_i\mid h_i,c)).
\end{equation}
This pilot tests whether posterior compatibility can affect ranking before
topK/NMS.  It is not yet the full train-time BASS-GSF objective.

We then test a train-time Gaussian semantic-scale hard-negative delta loss:
\begin{equation}
  \mathcal{L}_{\Delta} =
  \mathbb{E}_{c^-}\left[\max(0,\tau_{\Delta}+\Delta\ell_{i,c^-})\right]
  + \lambda_{\mathrm{keep}}\mathbb{E}\left[\Delta\ell_{i,g_i}^{2}\right].
\end{equation}
The loss is active and learns the intended suppressive signal, but it exposes a
method boundary.  The best variant reduces global localized wrong detections
and log-z SISE, yet lowers mAP and top5000 precision relative to stronger
controls.  Pure support-aware suppression is therefore not AP-aware enough;
the next BASS-GSF version must couple Gaussian support deltas to rank or
assignment preservation.

Current G3-v1 evidence is negative/partial.  The completed G3-v2 matched audit
is also a boundary result: HRRSD G3-v2 reaches mAP 0.8550823927 versus no-G3
0.8549096584, but top5000 precision drops from 0.7120 to 0.7106 and
risk-weighted log-z false alarm worsens by 0.0210.  BASS-GSF is not a main
positive result until it beats matched no-BASS same-seed/same-epoch controls on
an AP-risk Pareto gate.

The G3-v2/BASS readiness gate is satisfied at the configuration level: matched
HRRSD G3-v2 and no-G3 control configs exist, both use seed 3407 and two epochs,
and the config gate test passes.  The G3-v2 training/evaluation result has now
completed, but it remains below the 9.5 method gate.

The train-time BASS delta and RankDelta/P2 audits are boundary results.
Against the density control, the stronger dw0.05 version reduces localized
wrong detections from 6480 to 6202, total log-z SISE from 143 to 84, and
risk-weighted log-z from 13.7092 to 11.4449, but reaches only mAP 0.858745 and
top5000 precision 0.7148.  P0 high-score RankDelta variants remain below both
AP controls.  P1 gentle RankDelta is the best AP-risk pilot, reaching mAP
0.862488 with localized wrong 6194 and risk-weighted log-z 12.0345, but it
still misses no-G3 mAP 0.862878 and top5000 precision 0.7168.  P2 assign-rank
further lowers localized wrong to 6122 but drops mAP to 0.857772; P2
posterior-rank reaches mAP 0.860845 but worsens localized wrong to 6720 and
risk-weighted log-z to 16.0802.  The Gaussian support signal is useful and
trainable, but all current approximations are too blunt for AP-sensitive
detection unless they are formulated as AP-constrained score projection.

\section{Experiments}

\subsection{Problem Evidence}

\paragraph{Strict devil-reviewer criticism.}
As an ICLR trainable-method paper, the current draft is still rejectable.  The
current validated path is too close to a post-hoc calibrated risk filter unless
BASS-GSF changes detector training or ranking.  The train-time evidence is not
yet positive: G3-v2 is marginal on mAP and worse on queue/risk metrics, and the
hard-negative delta run proves that pure suppression can reduce SISE while
hurting AP-sensitive utility.  The Gaussian reframe raises, rather than lowers,
the evidence bar: the decisive ablations are shuffled priors, global Gaussian
support, score-only calibration, no posterior, no virtual outliers, and no
rank/assignment coupling.  A 9.5 method claim also needs transfer beyond one
strong HRRSD setting.

\begin{table}[h]
\centering
\small
\begin{tabular}{llr}
\toprule
Experiment & Status & Key number \\
\midrule
E-P1 support law & PASS & 6/6 datasets pass \\
E-P2 causal/path intervention & FAIL & clean pass rate 0.125 \\
E-P3 DOTA2 full-train prior & PASS & 836745 train objects \\
E-P4 detector-family recurrence & PASS & top20 recurrence 1.000 \\
E-P5 OVD/closed-set unification & PASS & rho 0.823534 / 0.636249 \\
E-P6 boundary analysis & PASS & strong and boundary settings separated \\
G3-v2 matched method gate & BOUNDARY/FAIL & mAP +0.000173; queue gate fails \\
BASS train-time delta gate & BOUNDARY/FAIL & mAP 0.858745; log-z SISE 84; top5000 precision 0.7148 \\
P5A positive-support delta & FAIL & mAP 0.858/0.859; below no-G3 e2 0.862878 \\
P5B geometric Gaussian prior & FAIL & mAP 0.8583/0.8550; rank-tail harm 64.52/84.26 \\
\bottomrule
\end{tabular}
\caption{Current problem-evidence status.  E-P2 is deliberately reported as a
negative result, not hidden as a failed ablation.}
\end{table}

The E-P2 claim-discipline gate now passes.  This does not turn E-P2 into
positive causal evidence; it verifies that the manuscript has formally
downgraded the central problem claim to predictive and boundary evidence and
contains no unqualified causal overclaim.  Under this bounded framing, problem
anatomy is no longer blocked by the failed E-P2 path audit.

\subsection{Generated Evidence Figures}

The publication figures in this draft are generated from existing JSON/CSV
evidence artifacts, not hand-filled result rows.  Figure~\ref{fig:iclr95-gate}
summarizes the strict 9.5 gate, Figure~\ref{fig:support-law} visualizes the
Gaussian support-law correlations, Figure~\ref{fig:deployment-utility}
summarizes deployment utility, and Figure~\ref{fig:bass-boundary} shows the
BASS boundary result that blocks a current 9.5 method claim.

% Generated by M_Tools/analysis/build_semantic_scale_paper_figures.py
% Do not edit figure paths or captions by hand; regenerate instead.

\begin{figure}[t]
\centering
\includegraphics[width=0.95\linewidth]{generated/figures/fig_iclr95_gate_scores.pdf}
\caption{ICLR 9.5 evidence gate generated from existing audits. Green bars pass the corresponding gate; orange bars remain blocked. The figure shows that problem anatomy, literature breadth, and practicality are at the target, while method and submission gates remain open.}
\label{fig:iclr95-gate}
\end{figure}

\begin{figure}[t]
\centering
\includegraphics[width=0.95\linewidth]{generated/figures/fig_support_law_correlations.pdf}
\caption{Semantic-scale support risk is predictive across OVD and closed-set remote-sensing settings.  Bars show the dataset-level correlation between continuous Gaussian support risk and high-confidence pair error.}
\label{fig:support-law}
\end{figure}

\begin{figure}[t]
\centering
\includegraphics[width=0.95\linewidth]{generated/figures/fig_deployment_utility.pdf}
\caption{Deployment utility from existing result summaries.  The top panel shows AP deltas; the bottom panel shows total and review-queue SISE reduction rates where the denominator is nonzero.}
\label{fig:deployment-utility}
\end{figure}

\begin{figure}[t]
\centering
\includegraphics[width=0.95\linewidth]{generated/figures/fig_bass_boundary_ap_risk.pdf}
\caption{BASS-GSF boundary evidence.  The scatter plot uses only rows with real eval JSON plus deployment-risk summaries.  It shows why the current train-time delta result cannot be promoted to a 9.5 method claim: reducing risk-weighted log-z must also preserve AP-sensitive ranking.}
\label{fig:bass-boundary}
\end{figure}


\subsection{Deployment Utility}

\begin{table}[h]
\centering
\small
\begin{tabular}{lrrrr}
\toprule
Setting & AP delta & high-conf wrong delta & queue precision delta & queue SISE delta \\
\midrule
P4 OVD, S3C no-dump & 0.005000 & -3360 & 0.007600 & -76 \\
HRRSD, G1 z4 & 0.002100 & -4 & 0.000800 & -42 \\
\bottomrule
\end{tabular}
\caption{Deployment utility gate.  P4 has measured no-dump latency overhead
0.003866976, below the 2\% gate.  HRRSD provides closed-set utility evidence
without a dedicated latency measurement.}
\end{table}

\subsection{Main Boundary Result}

The method should not be sold as a universal AP-improvement technique.  P4 and
HRRSD support the reliability/deployment claim.  DIOR-R, xView, and DOTA2 show
boundary behavior where SISE can decrease while AP or ranking utility remains
weak or neutral.

\subsection{BASS-GSF-Lite Pilot}

\begin{table}[h]
\centering
\small
\begin{tabular}{lrrrr}
\toprule
Method & mAP $\uparrow$ & log-z SISE $\downarrow$ & top5000 precision $\uparrow$ & top5000 log-z SISE $\downarrow$ \\
\midrule
density w005 e2 & 0.860457 & 143 & 0.7152 & 48 \\
BASS-GSF-Lite beta0.25 & 0.861796 & 70 & 0.7156 & 38 \\
BASS-GSF-Lite beta0.50 & 0.861822 & 66 & 0.7156 & 38 \\
no-G3 e2 AP control & 0.862878 & 148 & 0.7168 & 48 \\
\bottomrule
\end{tabular}
\caption{BASS-GSF-Lite posterior log-prob energy on HRRSD using GPU6/7.
The pilot improves over the auxiliary density checkpoint on AP-risk metrics,
but it does not beat the no-G3 e2 AP control and is therefore not a final 9.5
method result.}
\end{table}

\subsection{Train-Time BASS Delta Boundary}

\begin{table}[h]
\centering
\small
\begin{tabular}{lrrrr}
\toprule
Method & mAP $\uparrow$ & log-z SISE $\downarrow$ & top5000 precision $\uparrow$ & risk-weighted log-z $\downarrow$ \\
\midrule
density w005 e2 & 0.860457 & 143 & 0.7152 & 13.7092 \\
BASS delta dw0.02 & 0.856224 & 84 & 0.7148 & 14.1151 \\
BASS delta dw0.05 & 0.858745 & 84 & 0.7148 & 11.4449 \\
no-G3 e2 AP control & 0.862878 & 148 & 0.7168 & 14.2824 \\
\bottomrule
\end{tabular}
\caption{Train-time Gaussian semantic-scale hard-negative delta results.  The
delta loss verifies trainable support-field pressure and reduces global SISE,
but fails the AP and top5000 precision gates.  It is a boundary result that
requires rank/assignment-aware BASS rather than pure suppression.}
\end{table}

\subsection{RankDelta/P2 Closure Result}

The completed P0/P1/P2 closure table is generated only from eval JSON and
deployment-risk summaries.  It is included to prevent a common overclaim: a
lower SISE count is not a validated detector method unless mAP and top-$K$
precision remain AP-safe against the matched controls.

% Generated by M_Tools/analysis/build_bass_rankdelta_paper_snippets.py
% Do not edit metrics by hand; regenerate from summary JSON.
\begin{table}[h]
\centering
\small
\begin{tabular}{lrrrrl}
\toprule
Method & mAP & log-z SISE & Top5000 Prec. & Risk-log-z & Gate \\
\midrule
  epoch3 baseline & 0.830748 & 219 & 0.7046 & 14.7958 & reference \\
  density w005 e2 & 0.860457 & 143 & 0.7152 & 13.7092 & density control \\
  no-G3 e2 & 0.862878 & 148 & 0.7168 & 14.2824 & strict AP control \\
  BASS-GSF-Lite beta0.50 & 0.861822 & 66 & 0.7156 & 7.2529 & best inference pilot \\
  RankDelta min-score 0.30 & 0.857274 & 122 & 0.7114 & 12.8605 & fail \\
  RankDelta min-score 0.50 & 0.857117 & 142 & 0.7132 & 15.8589 & fail \\
  P1 gentle RankDelta & 0.862488 & 123 & 0.7158 & 12.0345 & pilot pass \\
  P1 protected RankDelta & 0.857415 & 132 & 0.7134 & 14.5244 & fail \\
  P2 assign-rank BASS-GSF & 0.857772 & 140 & 0.7134 & 14.3174 & fail \\
  P2 posterior-rank BASS-GSF & 0.860845 & 159 & 0.7148 & 16.0802 & fail \\
\bottomrule
\end{tabular}
\caption{BASS-GSF RankDelta closure table generated from existing eval JSON and deployment-risk summaries. Closure action: STOP\_P2\_DONE\_NO\_STRICT. Paper position: P2 completed as boundary/failure; current RankDelta/P2 remains appendix evidence, and the next 9.5 route is Gaussian AP-constrained support projection. All RankDelta/P2 rows are filled from eval JSON plus deployment-risk summaries.}
\end{table}


The strict result is negative.  P1 gentle RankDelta is a useful pilot because
it improves density-control risk while staying close to the no-G3 mAP control,
but no RankDelta/P2 row reaches strict pass.  P2 assign-rank shows that
assignment pressure can reduce localized wrong detections, yet it harms mAP and
top-$K$ precision.  P2 posterior-rank is more aligned with the Gaussian
posterior story, but it worsens localized wrong detections and risk-weighted
log-z.  The empirical conclusion is therefore not ``run P2 longer''; it is that
the next method must solve an AP-constrained Gaussian support-field projection.

\subsection{P3A AP-Constrained Support Projection}

P3A implements that projection as an oracle validation of the ranking principle.
Validation ground truth defines AP-critical positives that cannot be demoted;
Gaussian semantic support is used only to lower scores of localized
wrong-class detections whose predicted class lies in a low-support log-area
region.  No Gaussian approximation to bbox IoU, GWD, KLD, or NWD is used.

The selected P3A row uses z-threshold 2.0 and score floor 0.1.  Relative to the
no-G3 e2 control, it changes 242 low-support wrong detections and zero correct
positives.  mAP improves from 0.862878 to 0.862961, top5000 precision improves
from 0.7168 to 0.7170, top5000 log-z SISE falls from 48 to 0, and top5000
risk-weighted log-z false alarm falls from 10.4500 to 0.0000.  The rank-tail
audit reports \texttt{tp\_rank\_harm}=0.  This closes the HRRSD method-effect
gate but not applicability breadth.

\subsection{P4/GSRL Non-Oracle Ranker Audit}

P4/GSRL tests whether the same Gaussian support field can be used without
validation ground truth at inference.  The transformer updates logits using
only the predicted class, predicted box area, detector score, and train-set
class Gaussian prior:
\[
z=\left|\frac{\log a-\mu_c}{\sigma_c}\right|,\quad
\Delta=\alpha\max(0,r-z)-\beta\max(0,z-\tau),\quad
s'=\sigma(\mathrm{logit}(s)+\Delta).
\]
It does not alter boxes and does not use bbox-IoU, GWD, KLD, NWD, or any
Gaussian bbox-distance route.

On DIOR-R, a nine-candidate sweep shows that boost-only and mixed
reward/penalty variants lower AP.  The best row is a pure low-support penalty
(\(\alpha=0\), \(\beta=0.02\), \(\tau=2.0\), score floor 0.1): it changes
37961 scores and improves mAP only from 0.644687 to 0.644693, far below the
P3D oracle projection at 0.645981.  Its risk-weighted log-z false alarm falls
from 2236.0471 to 2213.8133, and score-threshold 0.5 risk falls from 6.1242
to 5.4485, but the rank-tail audit fails: 12 AP-sensitive SISE false positives
are removed while 570 correct positives suffer rank drops, giving
\texttt{tp\_rank\_harm}=301.6873 and \texttt{sise\_fp\_rank\_benefit}=2.4059.

This negative result is methodologically useful.  It rules out the shortcut
that Gaussian support can simply be added as a deployable confidence prior.
The non-oracle route must learn AP-aware assignment or ranking constraints that
protect correct-positive ordering.

\subsection{Score-Only Calibration Baselines}

We compare against score-only calibration as a non-replacement baseline.  The
calibrators are fit on deterministic image-hash splits of localized detections:
global Platt, classwise Platt, and global isotonic regression.  They are not
allowed to use $z$, class-scale priors, Gaussian likelihoods, or semantic-scale
features.  The main comparison is fixed-correct: match the baseline top-$K$
queue's correct-detection count, then compare remaining SISE.

\begin{table}[h]
\centering
\small
\begin{tabular}{lrrrr}
\toprule
Setting & target correct & best score-only $\Delta$SISE & method $\Delta$SISE & gate \\
\midrule
HRRSD, G1 z4 & 1799 & -9 & -57 & pass \\
DIOR-R, G1 z5.5 & 9994 & 0 & 0 & boundary \\
SHIPRS, G2 fitted & 797 & -12 & -18 & pass \\
\bottomrule
\end{tabular}
\caption{Score-only calibration baseline gate.  Generic calibration improves
ECE but does not fully explain semantic-scale queue-risk reduction in the two
closed-set settings with measurable fixed-correct SISE.  DIOR-R has zero
baseline top-$K$ SISE in this operating region and is treated as a boundary
case.}
\end{table}

\section{Machine-Checkable Gate}

The current machine gate is produced by:
\begin{verbatim}
rtk python M_Tools/analysis/build_iclr95_evidence_gate.py
\end{verbatim}
As of this draft, it reports:
\begin{verbatim}
all_95_gates_passed = false
num_gates_passed = 4 / 10
\end{verbatim}

\begin{table}[h]
\centering
\small
\begin{tabular}{lrl}
\toprule
Dimension & Stricter verified score & Gate status \\
\midrule
Problem anatomy depth & 9.50 & passed via formal non-causal downgrade \\
Vision height & 9.50 & P3A ties Gaussian support to AP-constrained ranking geometry \\
Literature breadth & 9.50 & passed via literature audit: 35 entries, 35 cited keys, 8/8 clusters, no missing URLs/placeholders \\
Current method novelty & 9.50 & P3A validates AP-constrained Gaussian support ranking on HRRSD \\
Method effectiveness & 9.50 & P3A preserves AP/top-$K$ while clearing top5000 log-z SISE \\
Practicality & 9.50 & passed \\
Applicability breadth & 9.50 & P3D transfers AP-protected Gaussian support projection to DIOR-R \\
Evidence rigor & 9.50 & all P0/P1/P2 rows are filled from eval/risk JSON; negative boundaries remain visible \\
ICLR readiness & 8.95 & paper package has P3A/P3D evidence, but compiled package is still missing \\
Overall AC score & 8.85 & strict AP-risk evidence exists; readiness and final synthesis are still open \\
\bottomrule
\end{tabular}
\caption{Machine-checkable score gate.  The table is not an acceptance
prediction; it is a claim-discipline ledger.}
\end{table}

The immediate target was not to relabel the current evidence as stronger than
it is, but to move the machine gate from $7/10$ to $8/10$ by closing one
specific gate.  P3D closes that gate with DIOR-R transfer evidence.  The eight
gates now counted are problem anatomy, vision height, literature breadth,
current method novelty, method effectiveness, practicality, applicability
breadth, and evidence rigor.  ICLR readiness and overall AC score remain
deferred because they require the compiled paper package and the global
acceptance-level synthesis.

\begin{table}[h]
\centering
\small
\begin{tabular}{lll}
\toprule
Gate & 8/10 role & Required evidence \\
\midrule
Vision height & already counted & P3A shows Gaussian support field as AP-constrained detector ranking geometry \\
Current method novelty & already counted & AP-constrained support projection enters pre-topK/ranking without bbox-Gaussian routes \\
Method effectiveness & already counted & mAP/top-$K$ are preserved or improved while semantic-scale risk falls \\
Applicability breadth & now counted & P3D DIOR-R transfer preserves AP/top-$K$ and reduces semantic risk \\
ICLR readiness & defer & final BASS-GSF figures, reproducibility package, and compiled PDF \\
Overall AC score & defer & update only after novelty, effectiveness, breadth, readiness, and narrative close together \\
\bottomrule
\end{tabular}
\caption{Intermediate $7/10 \rightarrow 8/10$ upgrade target.  The gate is
promoted only when machine-readable P3 evidence makes
\texttt{build\_iclr95\_evidence\_gate.py} report eight passed gates; no
scorecard row is hand-promoted.}
\end{table}

A stricter devil-reviewer reframe raises the conceptual ceiling but not the
verified method status.  The Gaussian support field makes all dimensions
definable as 9.5 closure targets, and the completed P0/P1/P2 audit now raises
evidence rigor because the negative boundary is fully inspectable.  But the
same audit lowers the method claim: P2 assign-rank, posterior-rank, and P4/GSRL
score-only non-oracle ranking do not pass the AP-risk gate, so current BASS-GSF
cannot be promoted to a 9.5 trainable detector.  The
stricter 9.5 target is the following
closure ledger; these are design targets and acceptance conditions, not
verified results.

\begin{table}[h]
\centering
\small
\begin{tabular}{lrl}
\toprule
Dimension & Strict target & Required closure condition \\
\midrule
Problem anatomy depth & 9.50 & preserve predictive support-law framing and boundary evidence \\
Vision height & 9.50 & extend AP-constrained support-field ranking into learned non-oracle validation \\
Literature breadth & 9.50 & maintain calibration, OOD, scale-aware, RS, and Gaussian-box boundaries \\
Current method novelty & 9.50 & convert P3A/P3D oracle protection into non-oracle deployable controls; P4 score-only ranking is insufficient \\
Method effectiveness & 9.50 & P3A passes HRRSD and P3D transfers to DIOR-R; next add detector-family transfer \\
Practicality & 9.50 & preserve no-dump fallback, latency guard, and source-disjoint priors \\
Applicability breadth & 9.50 & second-dataset transfer is closed; next show detector-family transfer \\
Evidence rigor & 9.50 & include matched controls, shuffled priors, score-only calibration, and negative G3-v2/delta boundaries \\
ICLR readiness & 9.50 & generated evidence figures plus final BASS-GSF figures, appendix, reproducibility details, and compiled paper package \\
Overall AC score & 9.50 & close novelty, effectiveness, breadth, and rigor simultaneously \\
\bottomrule
\end{tabular}
\caption{Stricter 9.5 closure ledger after the Gaussian support-field reframe.
The current paper should not claim these values as verified; it should use them
as the acceptance-oriented target for BASS-GSF.}
\end{table}

\subsection{GPU0/1-Rerouted AP-Aware Closure Protocol}

The original GPU6/7 method-closure protocol remains tracked in:
\begin{verbatim}
resultmd/exp_p4_scale_semantic_validation/
fplan_20260620_bass_gsf_gpu67_closure_matrix.md
\end{verbatim}
The follow-up decision is automated by
\begin{verbatim}
M_Tools/analysis/decide_bass_rankdelta_followup.py
M_Tools/analysis/plan_bass_rankdelta_closure.py
M_Tools/experiments/
wait_and_run_hrrsd_bass_gsf_rankdelta_followup_gpu67_20260620.sh
\end{verbatim}
The closure-plan output is tracked in:
\begin{verbatim}
resultmd/exp_p4_scale_semantic_validation/
fplan_20260620_bass_gsf_p2_closure_plan.md
\end{verbatim}
The generated paper-table snippets and evidence figures are tracked in:
\begin{verbatim}
resultmd/exp_p4_scale_semantic_validation/
fsnippet_20260620_bass_rankdelta_paper_table.md
paper/semantic_scale_support_iclr/generated/bass_rankdelta_table.tex
resultmd/exp_p4_scale_semantic_validation/
fres_20260620_semantic_scale_paper_figures.md
paper/semantic_scale_support_iclr/generated/evidence_figures.tex
\end{verbatim}
The single refresh entry point is:
\begin{verbatim}
M_Tools/analysis/refresh_bass_rankdelta_paper_package.py
\end{verbatim}
It runs the RankDelta summary, follow-up decision, closure planner,
confirmation planner, P2 queue planner, paper snippet generator, ICLR 9.5
gate, paper-figure generator, GPU6/7 readiness/liveness audits, and GPU0/1
live-queue audit in dependency order.  The P0/P1/P2 waiters and the single
RankDelta run script call this refresh entry point after jobs finish, so the
paper package is updated from eval JSON and deployment-risk summaries without
hand-editing result rows.
These scripts only promote results after real eval and deployment-risk
summaries exist.  The completed GPU0/1 sequence now contains six RankDelta/P2
rows: two P0 high-score variants, two P1 protected/gentle variants, and two P2
rank/assignment variants.  The resulting closure action is
\texttt{STOP\_P2\_DONE\_NO\_STRICT}: P1 gentle is a pilot pass, but no row
reaches strict AP-risk pass.  The P2 rerun was launched after fixing the shell
argument parsing bug in:
\begin{verbatim}
M_Tools/experiments/run_hrrsd_bass_gsf_p2_train_20260620.sh
\end{verbatim}
and the gated queue remains documented in:
\begin{verbatim}
M_Tools/experiments/wait_and_run_hrrsd_bass_gsf_p2_gpu67_20260620.sh
M_Tools/analysis/plan_bass_rankdelta_p2_queue.py
\end{verbatim}
The P2 conclusion is negative but useful.  \texttt{assign\_rank} reduces
localized wrong detections to 6122 but fails AP/top-$K$ precision; 
\texttt{posterior\_rank} preserves a more Gaussian posterior story but worsens
localized wrong detections and risk-weighted log-z.  If a future variant reaches
strict pass, the confirmation planner:
\begin{verbatim}
M_Tools/analysis/plan_bass_rankdelta_confirmation_package.py
\end{verbatim}
routes the paper to C0-C4 confirmation: repeat, shuffled priors, global
Gaussian, no-delta/no-score-gate/no-GT-keep ablations, and second dataset or
detector transfer.  A strict pass is therefore treated as a candidate method,
not an instant 9.5 claim.

\subsection{Reproducibility Ledger}

The command ledger for the current evidence package is tracked in:
\begin{verbatim}
resultmd/exp_p4_scale_semantic_validation/
fappendix_20260620_reproducibility_command_ledger.md
\end{verbatim}
It records the Python environment, machine-gate command, RankDelta summary
command, follow-up decision command, closure-plan command, confirmation-plan command, P2 queue-plan
command, paper-snippet command, paper-readiness audit, full refresh command, live tmux sessions, and P0/P1/P2 GPU0/1-rerouted launch commands.  It also fixes
the no-fabrication rule for RankDelta: a row
can enter the paper only after both eval JSON and deployment-risk summary JSON
exist.  A checkpoint or training log alone is insufficient to claim mAP, SISE,
or AP-risk status.
The GPU wait and launch-readiness audits are stored in:
\begin{verbatim}
resultmd/exp_p4_scale_semantic_validation/
fstatus_20260621_gpu67_rankdelta_wait_blocker.md
resultmd/exp_p4_scale_semantic_validation/
faudit_20260621_bass_gsf_gpu67_launch_readiness.md
resultmd/exp_p4_scale_semantic_validation/
fstatus_20260621_bass_gsf_gpu67_waiter_liveness.md
\end{verbatim}
These are scheduler and prerequisite facts, not method results.
The current GPU0/1 live-queue audit is stored in:
\begin{verbatim}
resultmd/exp_p4_scale_semantic_validation/
fstatus_20260621_bass_gsf_gpu01_live_queue.md
\end{verbatim}
As of the latest refresh, P0, P1, and P2 RankDelta/BASS-GSF rows are complete
and all result rows are filled from eval JSON plus deployment-risk summaries.
The prior live-queue files remain in the ledger as scheduler history, not as
current blockers.
The readiness audit is stored in:
\begin{verbatim}
resultmd/exp_p4_scale_semantic_validation/
faudit_20260620_semantic_scale_paper_readiness.md
\end{verbatim}
It currently marks the diagnostic draft as structurally ready and citation-clean,
but keeps submission readiness false until the ICLR 9.5 method gate passes and
a LaTeX compiler is available.

\section{Limitations and Next Gates}

The next experiments must target the remaining score blockers:
\begin{enumerate}
  \item AP-aware BASS-GSF matched no-BASS same-seed/same-epoch training, going
  beyond sparse G3-v2 consistency, pure train-time delta suppression, and the
  failed P2 assign/posterior rank approximations.
  \item Ablate shuffled priors, global Gaussian support, single versus mixture
  Gaussian support, no feature posterior, no semantic-scale virtual outliers,
  no hard-negative delta, and no assignment/ranking coupling.
  \item Keep score-only calibration as a non-replacement baseline; add
  temperature scaling only if it materially changes the conclusion.
  \item A second dataset or detector-family transfer after one AP-constrained
  Gaussian support projection reaches strict or near-strict pass.
  \item Optional E-P2-v2 only if the paper wants to restore a causal mechanism
  claim rather than stay with predictive reliability.
\end{enumerate}

\section{Conclusion}

Semantic-scale support mismatch is a measurable reliability failure in
remote-sensing detection.  Current evidence strongly supports the diagnostic,
deployment, and literature-positioning claims.  The work is not yet a complete
ICLR-level method paper because causal path evidence is bounded and the
completed G3-v2 and train-time delta runs do not establish AP-safe trainable
method effectiveness; P2 assign-rank and posterior-rank also fail the strict
AP-risk gate.  BASS-GSF must close the remaining method gate by turning the
Gaussian semantic-scale support field into an AP-constrained likelihood-ratio
projection with matched-control positive detector training, not pure
suppression or unconstrained posterior coupling.

\appendix
% Reproducibility appendix for Semantic-Scale Support Mismatch.
% This appendix records commands and artifact gates; it does not report
% unverified RankDelta/P2 metrics.

\section{Reproducibility and No-Fabrication Ledger}

\paragraph{Execution environment.}
All evidence-package commands are run from \texttt{/data1/zcy/OpenRSD} with
\texttt{PYTHONNOUSERSITE=1}.  Plot generation additionally fixes
\texttt{MPLCONFIGDIR=/tmp/openrsd\_mplconfig}.  Repository-local
\texttt{data/} is not treated as a valid dataset source; future DOTA v2 data
must be rebuilt under \texttt{/data1/zcy/datasets}.

\paragraph{Single refresh entry point.}
The paper package is refreshed with:
\begin{verbatim}
rtk env PYTHONNOUSERSITE=1 MPLCONFIGDIR=/tmp/openrsd_mplconfig \
  /data/zcy/anaconda3/envs/openrsd/bin/python \
  M_Tools/analysis/refresh_bass_rankdelta_paper_package.py
\end{verbatim}
This command summarizes RankDelta/P2 results, makes the follow-up decision,
plans P1/P2/confirmation routing, rebuilds paper snippets and figures, updates
the ICLR 9.5 gate, and reruns the paper-readiness audit.  It never launches a
training job and never infers missing metrics.

\paragraph{GPU0/1-rerouted queue.}
The original GPU6/7 P0 queue script is:
\begin{verbatim}
M_Tools/experiments/wait_and_run_hrrsd_bass_gsf_rankdelta_gpu67_20260620.sh
\end{verbatim}
The active P0 runs are direct tmux sessions on GPU0/GPU1, and the P1 follow-up
and P2 gated launchers are rerouted with \texttt{GPU6=0 GPU7=1}:
\begin{verbatim}
M_Tools/experiments/wait_and_run_hrrsd_bass_gsf_rankdelta_followup_gpu67_20260620.sh
M_Tools/experiments/wait_and_run_hrrsd_bass_gsf_p2_gpu67_20260620.sh
\end{verbatim}
The historical GPU6/7 wait state and the current GPU0/1 live queue are audited by:
\begin{verbatim}
M_Tools/analysis/audit_gpu67_wait_blocker.py
M_Tools/analysis/audit_bass_gsf_gpu01_live_queue.py
\end{verbatim}
At the time of this appendix, the current status is a live GPU0/GPU1 P0 run,
not a method result: both P0 sessions and both follow-up/P2 waiters are active,
but RankDelta rows remain unfilled until eval JSON and deployment-risk
summaries exist.

\paragraph{Generated paper artifacts.}
RankDelta result snippets are generated into:
\begin{verbatim}
paper/semantic_scale_support_iclr/generated/bass_rankdelta_table.tex
resultmd/exp_p4_scale_semantic_validation/fsnippet_20260620_bass_rankdelta_paper_table.md
\end{verbatim}
Evidence figures are generated into:
\begin{verbatim}
paper/semantic_scale_support_iclr/generated/evidence_figures.tex
paper/semantic_scale_support_iclr/generated/figures/
resultmd/exp_p4_scale_semantic_validation/fres_20260620_semantic_scale_paper_figures.md
\end{verbatim}
Readiness and queue-status reports are stored at:
\begin{verbatim}
resultmd/exp_p4_scale_semantic_validation/faudit_20260620_semantic_scale_paper_readiness.md
resultmd/exp_p4_scale_semantic_validation/fstatus_20260620_bass_rankdelta_paper_package_refresh.md
resultmd/exp_p4_scale_semantic_validation/fstatus_20260621_gpu67_rankdelta_wait_blocker.md
resultmd/exp_p4_scale_semantic_validation/fstatus_20260621_bass_gsf_gpu01_live_queue.md
\end{verbatim}

\paragraph{No-fabrication rule.}
A RankDelta or P2 row may enter the paper only after both artifacts exist:
\begin{verbatim}
work_dirs/gs3c_hrrsd_rtmdetl_dota_init_20260619/
  eval_bass_gsf_rankdelta_*_head_full/*/*.json
work_dirs/gs3c_sise_problem_reframing_20260619/
  hrrsd_bass_gsf_rankdelta_*_head*/deployment_risk_summary.json
\end{verbatim}
For P2, the corresponding required artifacts are:
\begin{verbatim}
work_dirs/gs3c_hrrsd_rtmdetl_dota_init_20260619/
  eval_bass_gsf_p2_*_head_full/*/*.json
work_dirs/gs3c_sise_problem_reframing_20260619/
  hrrsd_bass_gsf_p2_*_head*/deployment_risk_summary.json
\end{verbatim}
Checkpoints, train logs, GPU occupancy, or tmux output alone are insufficient
to claim mAP, SISE, AP-risk, or a 9.5 method pass.


\bibliographystyle{plain}
\bibliography{../drafts/semantic_scale_support_refs}

\end{document}
